\chapter{Summary and outlook}
\label{conclusion}

In summary, in this thesis we have explored the application of boosted decision trees to the separation of $D^0$ signals and background at the EIC. In Ch. \ref{ch:d0_top}, we examined the features of $D^0\rightarrow\pi+K$ decay topology which we utilize to distinguish true $D^0$ signals from background, and we noted the dependence of the topological distributions on $y$ and $p_T$ of the $D^0$. In Ch. \ref{ch.ml}, we introduced the BDT methodology which we applied to $D^0$ analysis in Ch. \ref{ch.results}. We first utilized an existing framework \cite{Shyam_ml} to analyze new simulated data samples in Ch. \ref{sec.results_bkg_study}, producing a preliminary analysis which demonstrated the negative impacts of machine-background effects on the effectiveness of $D^0$ invariant mass reconstruction. We then improve upon the BDT framework in Ch. \ref{sec.results_optimization}, separately analyzing the data in slices of $y$ and $p_T$, optimizing the BDT threshold to maximize the Significance, testing the various combinations of topological features, and adjusting the selection of BDT hyperparameters using Optuna. 

Through these analyses, we observed varying results in different Ch. \ref{sec.results_prelim} slices, with the forward $y$ and low $p_T$ region exhibiting the least effective signal-and-background separation but the greatest improvement in Significance after BDT application. In Ch. \ref{sec.results_features}, we found that changing the feature selection can improve the results in some kinematic slices, particularly ones with lower statistics, while producing negligible effect in others. We note that more features does not generally lead to better BDT performance. Then in Ch. \ref{sec.results_hyperparams}, we observed a tendency for Optuna to select hyperparameter values which lead to overfitting, suggesting the need to restrict the ranges of basic hyperparameters to more conservative values. Finally, in Ch. \ref{sec.results_box_cuts} we illustrated the greater flexibility and greater improvements provided by BDT based analysis compared to traditional box cuts on topological features.

In future work, the steps taken in Ch. \ref{sec.results_optimization} to optimize the BDT framework can be applied to improve the analysis of data samples with machine-background effects, which would provide a more realistic demonstration of $D^0$ invariant mass reconstruction for actual EIC experiments. Additionally, there remains multiple avenues to explore in the tuning of BDT hyperparameters with Optuna. Firstly, expanding the range of regulatory hyperparameters and/or optimization of additional hyperparemeters (listed in Ch. \ref{sec.ml_optuna}) may be able to effectively prevent overfitting. Secondly, Optuna provides different sampling and optimization algorithms \cite{Optuna_docs} which it uses to traverse the parameter space, which can also impact the resulting hyperparameter values selected. These aspects can be examined to better understand and control the tendency of overfitting for different kinematic slices.